\begin{abstract}
Deploying Vision-Language-Action (VLA) models for real-time control requires compressing their vision token sequences, yet identical pruning applied to different architectures produces starkly divergent outcomes, from unscathed performance to total collapse. We trace this divergence to the vision-language connector, the module bridging the visual encoder and language backbone. Across six VLA architectures (450\,M--7\,B parameters) and four LIBERO manipulation suites, a single forward-pass diagnostic, the coefficient of variation (CV) of post-connector token norms, predicts pruning tolerance without any task evaluation. Low-CV connectors permit aggressive pruning; high-CV connectors collapse when even a few tokens are removed. No scoring method is universally safe: the advantage of importance-based over random selection follows an inverted-U curve that peaks at moderate compression and reverses at aggressive ratios, with task complexity amplifying connector-dependent patterns. A zero-mask ablation reveals that position disruption, not information loss, is the binding constraint. For architectures where all deletion fails, a lightweight re-encoder (${\sim}$5\% extra parameters) recovers full task performance at 75\% compression.
\end{abstract}
